Let $\mathbf{x} = \{x[n]\}_{n \in \mathbb{Z}}$ be a discrete-time signal representing the input time series. We define the patch-based tokenization operator $\mathcal{T}_{P,S}$ with patch length $P \in \mathbb{N}^{+}$ and stride $S \in \mathbb{N}^{+}$ as a mapping that transforms the signal $\mathbf{x}$ into a sequence of vector-valued tokens $\mathbf{T}=\{\mathbf{t}_k\}_{k\in\mathbb{Z}}$, where each token $\mathbf{t}_k\in\mathbb{R}^{P}$.

The $k$-th token is defined as

\begin{equation}
\label{eq:token_def}
\mathbf{t}_k
=
\mathcal{T}_{P,S}(\mathbf{x})[k]
=
\begin{bmatrix}
x[kS] \\
x[kS+1] \\
\vdots \\
x[kS+P-1]
\end{bmatrix}.
\end{equation}

Using element-wise indexing notation,

\begin{equation}
\label{eq:element_def}
\mathbf{t}_k[m]
=
x[kS+m],
\qquad
m\in\{0,\ldots,P-1\}.
\end{equation}

Before examining the degeneracies of this operator it is worth fixing what it does
not do. Whenever $S\le P$, every sample of the input appears in at least one
token and can be read back from it: taking $k=\lfloor n/S\rfloor$ and
$m = n \bmod S \le S-1 \le P-1$ in \autoref{eq:element_def} gives

\begin{equation}
\label{eq:reconstruction}
x[n]
=
\mathbf{t}_{\lfloor n/S\rfloor}\!\left[n \bmod S\right],
\qquad
S\le P.
\end{equation}

Complete raw patching with $S\le P$ is therefore injective: the token sequence
determines the signal, and no frequency below Nyquist is destroyed by tokenization
itself. Everything that follows concerns redundancy in the token sequence
rather than loss of information, the single exception being the $S>P$ regime of
Case 3 below, where the operator does discard samples.